\documentclass{article}
\usepackage{xcolor}
\usepackage{amssymb}
\usepackage{amsmath}

\title{Quiz 12}
\author{Brandon Reich}
\date{21 April 2026}

\begin{document}
\maketitle

\subsection*{Question 1}
In section 9.14, the textbook states Pascal's identity: for positive $n$ and $k$ such that $k < n$, $\binom{n}{k} = \binom{n-1}{k-1} + \binom{n-1}{k}$. The following proof for Pascal's identity was generated by AI. Is the proof correct? If so, explain why. If not, explain what parts of the proof are incorrect. \\

\textcolor{blue}{
The proof is not correct. There are a couple of problems with it.
} \\

\textcolor{blue}{
First, the base case doesn't work. The identity is only stated for $k < n$, but the base case uses $n = k$, which breaks that condition. Since the induction is on $n$ with $k$ fixed, the smallest value of $n$ we actually care about is $n = k + 1$, so that should be the base case.
} \\

\textcolor{blue}{
The bigger problem is the inductive step. The whole point of the inductive step is to use the inductive hypothesis to show the identity holds for $n = m+1$. But the proof doesn't really do that. It just writes out the identity for $n = m+1$ and says it's true because of Pascal's identity, which is the thing we're trying to prove in the first place. That's circular. Also the inductive hypothesis is never actually used anywhere, so there isn't really an induction happening.
} \\

\textcolor{blue}{
So even though Pascal's identity is true, this proof doesn't actually show it. The base case is outside the domain the textbook gave, and the inductive step assumes what it's trying to prove.
}

\subsection*{Question 2}
Prove that $\sum_{k=0}^{n} k\binom{n}{k} = n 2^{n-1}$, $n \geq 0$. Hint: recall that $\sum_{k=0}^{n} \binom{n}{k} = 2^n$, $n \geq 0$. \\[0.5em]

\textcolor{blue}{\textbf{Proof:}} \\[0.3em]

\textcolor{blue}{
First I want to use the identity $k\binom{n}{k} = n\binom{n-1}{k-1}$. To show this, just expand both sides using the factorial definition:
\[
k\binom{n}{k} = k \cdot \frac{n!}{k!(n-k)!} = \frac{n!}{(k-1)!(n-k)!},
\]
\[
n\binom{n-1}{k-1} = n \cdot \frac{(n-1)!}{(k-1)!(n-k)!} = \frac{n!}{(k-1)!(n-k)!}.
\]
Both sides match, so the identity holds.
} \\[0.5em]

\textcolor{blue}{
For $n = 0$, the left side is $0 \cdot \binom{0}{0} = 0$ and the right side is $0 \cdot 2^{-1} = 0$, so the formula works.
} \\[0.5em]

\textcolor{blue}{
For $n \geq 1$, start with the sum and drop the $k = 0$ term since it's zero:
\[
\sum_{k=0}^{n} k\binom{n}{k} = \sum_{k=1}^{n} k\binom{n}{k}.
\]
Now apply the identity from above:
\[
\sum_{k=1}^{n} k\binom{n}{k} = \sum_{k=1}^{n} n\binom{n-1}{k-1} = n \sum_{k=1}^{n} \binom{n-1}{k-1}.
\]
Reindex with $j = k - 1$ so the sum goes from $j = 0$ to $j = n-1$:
\[
n \sum_{k=1}^{n} \binom{n-1}{k-1} = n \sum_{j=0}^{n-1} \binom{n-1}{j}.
\]
By the hint, $\sum_{j=0}^{n-1} \binom{n-1}{j} = 2^{n-1}$, so
\[
n \sum_{j=0}^{n-1} \binom{n-1}{j} = n \cdot 2^{n-1}.
\]
That gives $\sum_{k=0}^{n} k\binom{n}{k} = n \cdot 2^{n-1}$ for all $n \geq 0$. $\blacksquare$
}

\end{document}
